\section{The \texttt{Mirror} McXtrace Component}
Perfectly flat mirror (in XZ or YZ), or polygonal

\subsection*{Identification}
\begin{itemize}
  \item \textbf{Author:} Erik B Knudsen
  \item \textbf{Origin:} DTU Physics
  \item \textbf{Date:} July 2016
\end{itemize}

\subsection*{Description}
This is a simple implementation of a perfectly flat mirror. The mirror plane is in the XZ-plane. It can be oriented in the YZ plane by setting 'yheight'. It may also be a complex polygonal geometry (OFF/PLY) by setting 'geometry'.

Reflectivity may be specified either as a number (R0) or by means of a material datafile. The material datafile may be specified as a coating or as relfectivity - either parameterized by q or E,theta. If the datafile is identified as a coating recipe, an ab-initio reflectivity calculation is triggered.

Example: Mirror(xwidth=5e-2, zdepth=2e-1, R0=1, coating="B4C.dat")

\subsection*{Input parameters}
Parameters in \textbf{boldface} are required; the others are optional.

\begin{longtable}{p{0.22\textwidth}p{0.12\textwidth}p{0.46\textwidth}p{0.14\textwidth}}
\toprule
\textbf{Name} & \textbf{Unit} & \textbf{Description} & \textbf{Default} \\
\midrule
\endhead
zdepth & m & The length of the mirror & 0.1 \\
xwidth & m & The width of the mirror & 0.01 \\
yheight & m & The height of the mirror. This overrides xwidth and puts the mirror in the yz-plane. & 0 \\
coating & str & Filename containing reflectivities (or coating). & "" \\
R0 & 0-1 & Constant reflectivity & 0 \\
geometry & str & Filename of an OFF/PLY geometry providing a polygonal surface. When xwidth/yheight/zdepth are also given, the object is rescaled accordingly. & "" \\
\bottomrule
\end{longtable}

\subsection*{Links}
\begin{itemize}
  \item Component source code found in file \texttt{Mirror.comp}.
\end{itemize}
\IfFileExists{optics/Mirror_static.tex}{\input{optics/Mirror_static.tex}}{}